\documentclass[10pt,a4paper]{article}
\usepackage[T1]{fontenc}
\usepackage[utf8]{inputenc}
\usepackage{lmodern}
\usepackage{amsmath,amssymb}
\usepackage[margin=21mm]{geometry}
\usepackage{microtype}
\usepackage[hidelinks]{hyperref}
\usepackage{enumitem}
\setlist[itemize]{leftmargin=*,nosep,topsep=3pt}
\setlength{\parindent}{0pt}
\setlength{\parskip}{5pt}
\setlength{\emergencystretch}{3em}
\newcommand{\question}[1]{\par\medskip{\normalsize\bfseries #1}\par\smallskip}
\pagestyle{plain}

\begin{document}
{\Large\bfseries Insights and opportunities for deep brain stimulation as a brain circuit intervention\par}

\textbf{Paper:} Neumann, Horn and Kuehn (2023).\quad Review.\par
\href{https://doi.org/10.1016/j.tins.2023.03.009}{DOI: \nolinkurl{10.1016/j.tins.2023.03.009}}\par
{\footnotesize Page references below use PDF page numbers, including front matter.}

\question{To which general question does this paper relate?}
Measurement, stimulation targets, symptom mechanisms and feedback objectives. Connects local physiology with whole-network DBS effects.

\question{How does it answer the general questions?}
\textbf{Measured and controlled.} The review connects unit activity, LFP oscillations, cortical signals and imaging with stimulation of neural circuits. DBS can change synchrony and symptoms through distributed pathways; the effect depends on anatomy and stimulation settings (PDF pp.~1--8).

\textbf{Usable insight.} Connectomic targeting and physiological feedback address complementary spatial and temporal aspects. Adaptive stimulation can use symptom-linked biomarkers, but the review explicitly leaves their causal role open. Physiological and pathological synchrony must be distinguished before selecting an objective (PDF pp.~5--6, 12).

\textbf{Model relevance.} Supports choosing measurement and actuation channels around an explicit symptom/network hypothesis. Tract correlations alone cannot establish whether a pathway is activated or inhibited, or the direction of propagation. The review supplies neither a single fitted delayed rate model nor a PIE multiplier construction (PDF pp.~8--12).

\question{What papers does it cite that could provide more information?}
\begin{itemize}
\item Little et al. (2013), \emph{Adaptive deep brain stimulation in advanced Parkinson disease}: early biomarker-driven intervention (ref.~73, PDF p.~14).
\item Gradinaru et al. (2009), \emph{Optical deconstruction of parkinsonian neural circuitry}: pathway-specific causal interrogation (ref.~147, PDF p.~16).
\end{itemize}

\end{document}
